% Part: normal-modal-logic
% Chapter: axioms-systems
% Section: soundness

\documentclass[../../../include/open-logic-section]{subfiles}

\begin{document}

\olfileid[fr]{nml}{prf}{snd}

\olsection{Correction}

Un système de !!{derivation} est dit correct si toute !!{formula}
qu'on peut y !!{derive} est valide. Pour les systèmes modaux, où l'on
peut employer, outre \Ax{K}, des instances des !!{formula}s
$!A_1$, \dots,~$!A_n$, on souhaite que chaque !!{formula}
!!{derivable} soit vraie dans tout modèle où toutes les instances de
$!A_1$, \dots, $!A_n$ sont vraies.

\begin{thm}[Théorème de correction]\ollabel{thm:soundness}
  Si toutes les instances de $!A_1$, \dots, $!A_n$ sont valides,
  respectivement, dans les classes de modèles $\mClass{C}_1$, \dots,
  $\mClass{C}_n$, alors $\Log{K}!A_1\dots !A_n \Proves !B$ implique
  que $!B$ est valide dans la classe de modèles
  $\mClass{C}_1 \cap \dots \cap \mClass{C}_n$.
\end{thm}

\begin{proof}
  Raisonnons par récurrence sur la longueur des démonstrations. Pour
  abréger, posons $\mClass{C} = \mClass{C}_1 \cap \dots \cap
  \mClass{C}_n$.
  \begin{enumerate}
  \item Cas de base : si $!B$ possède une démonstration de
    longueur~$1$, alors il s'agit soit d'une instance de tautologie, soit
    d'une instance de~\Ax{K}\iftag{prvDiamond}{ ou de~\Dual{}}{}, soit
    d'une instance de l'une des !!{formula}s $!A_1$, \dots,~$!A_n$.
    Dans le premier cas, $!B$ est valide dans $\mClass{C}$, puisque
    les instances de tautologies sont valides dans \emph{toute}
    classe de modèles, d'après la \olref[syn][tau]{prop:valid-taut}.
    Il en va de même dans le deuxième cas, d'après la
    \olref[syn][sch]{prop:Kvalid}\iftag{prvDiamond}{ et la
      \olref[syn][sch]{prop:Dual-valid}}{}. Enfin, dans le troisième
    cas, $!B$ est valide dans $\mClass{C}_i$ et
    $\mClass{C} \subseteq \mClass{C}_i$ ; donc $!B$ est également
    valide dans $\mClass{C}$ d'après la
    \olref[syn][val]{prop:subset-class}.
  \item Étape de récurrence : supposons que $!B$ possède une
    démonstration de longueur $k>1$. Si $!B$ est une instance de
    tautologie, de~\Ax{K}\iftag{prvDiamond}{ ou de~\Dual{}}{}, ou
    encore de l'une des !!{formula}s $!A_1$, \dots, $!A_n$, on
    raisonne comme dans le cas de base. Supposons donc que $!B$ soit
    obtenu par \MP{} à partir des !!{formula}s précédentes
    $!C \lif !B$ et $!C$. Alors $!C \lif !B$ et $!C$ possèdent des
    démonstrations de longueur $<k$ et sont, par hypothèse de
    récurrence, valides dans
    $\mClass{C}$. D'après la \olref[syn][sch]{prop:soundMP}, $!B$ est
    également valide dans $\mClass{C}$. Supposons enfin que $!B$ soit obtenu par
    \Nec{} à partir de $!C$, de sorte que $!B = \Box!C$. Par
    hypothèse de récurrence, $!C$ est valide dans $\mClass{C}$.
    Dans chaque modèle de cette classe, $!C$ est donc vraie en
    tous les successeurs de chaque monde. Le raisonnement de la
    \olref[syn][val]{prop:Nec-rule} donne ainsi la validité de~$!B$
    dans~$\mClass{C}$.
  \end{enumerate}
\end{proof}
\end{document}
